\documentclass[UTF8,fontset=fandol,10pt]{ctexart}
\usepackage[a4paper,margin=18mm,top=20mm,bottom=19mm]{geometry}
\usepackage[most]{tcolorbox}
\usepackage{enumitem,fancyhdr}
\definecolor{Navy}{HTML}{245781}
\definecolor{Risk}{HTML}{AC161C}
\definecolor{Good}{HTML}{167333}
\definecolor{Gold}{HTML}{9B8B39}
\definecolor{Slate}{HTML}{526273}
\pagestyle{fancy}\fancyhf{}
\setlength{\headheight}{15pt}
\fancyfoot[C]{\small\thepage}
\setlength{\parindent}{0pt}\setlength{\parskip}{4pt}
\setlist[itemize]{leftmargin=1.2em,itemsep=2pt,topsep=1pt}
\ctexset{section={format=\large\bfseries,beforeskip=12pt,afterskip=7pt},subsection={format=\normalsize\bfseries}}
\raggedbottom
\fancyhead[L]{\small 虚构案例：内容运营面试复盘}\fancyhead[R]{\small 问题逐字稿}
\begin{document}\begin{center}{\large 问题逐字稿}\par\vspace{4mm}{\LARGE\bfseries 虚构案例：内容运营面试复盘}\par\vspace{2mm}{\large\bfseries 问题与现场回答原话}\par\vspace{3mm}{\small 虚构演示 · 内容运营 · 无时间戳文字稿}\end{center}\textbf{版本与证据边界：}本例来自虚构文字稿；无录音，不代表已经进行人工听校。\par 全部材料均为虚构演示。三类材料齐全，但无时间戳，第三题回答缺失；仅评价记录支持的内容，不预测录用结果。\par\textbf{阅读说明：}保留口头语、自我更正、追问及不确定标记；无时间戳时使用段落编号。\par\section*{活动经历与个人职责}\textbf{[P1—P2] 面试官}\par 请介绍你的活动经历。\par \textbf{候选人}\par 我参与过一次内容活动，主要整理素材、跟进发布。具体业绩我没有数据。\par \section*{播放量下降的排查思路}\textbf{[P3—P4] 面试官}\par 如果播放量下降，你会如何排查？\par \textbf{候选人}\par 我会先比较同一账号最近几条内容，再看是不是题材或者发布时间不一样。\par \section*{跨团队协作与分歧处理}\textbf{[P5—P6] 面试官}\par 你怎样处理和设计同事的分歧？\par \textbf{候选人}\par [该段缺失，无法核对回答]\par \end{document}